% =============================================================================
\section{Physics of the memory chain}\label{sec:statmech}

\subsection{The stationary law}\label{sec:stationary}

With one lag, $\theta(x)=h+Jx$, the causal chain~\eqref{eq:causal} satisfies detailed balance with respect to
\begin{equation}
\pi(x)\propto e^{hx}\,Z(h+Jx),\qquad \pi(x)\,p(y\,|\,x)\propto e^{h(x+y)+Jxy},
\label{eq:stationary}
\end{equation}
as one checks by inserting Eq.~\eqref{eq:single} (Appendix~\ref{app:single}). The joint law of two consecutive flights is a two-site Boltzmann weight of Ising type, with continuous spins on $[0,L]$, and every stationary property of the walk follows from it.

The stationary law has a simple reading. In a reversible chain, the weight of a flight $x$ is the total weight of the pairs it forms with its successor, $\int_0^Ldy\,e^{h(x+y)+Jxy}=e^{hx}Z(h+Jx)$. The first factor is the bare power law. The second counts the successors: for $J>0$ a long flight forms favorable pairs with long successors, so the factor $Z(h+Jx)$ grows with $x$, roughly as $e^{JxL}$ once the field $h+Jx$ is positive, and it lifts the tail of the distribution. Figure~\ref{fig:overview}(d) shows the effect at the three states marked in the phase diagram of \figref{fig:overview}(c), all with the same bare exponent. In A, a weak coupling, the stationary law is barely distinguishable from the renewal one; in B the tail bends upward before the cutoff; in C most flights sit at $\lmax$.

\subsection{Duality and condensation}\label{sec:condensation}

Exchange long and short flights, $\ell\to\lmin\lmax/\ell$, that is, $x\to L-x$. For a renewal walker this maps the exponent $\mu$ onto $2-\mu$, a law that favors short flights onto one that favors long ones. With memory, the pair weight $e^{h(x+y)+Jxy}$ keeps its form: the coupling is unchanged and the field becomes $-h-JL$, up to a constant. A chain at $(h,J)$ and one at $(-h-JL,J)$ are thus the same chain read upside down, and the inversion is a duality. The chains with $h=-JL/2$ are mapped onto themselves. Their stationary law is symmetric under inversion, so exactly half of their flights are longer than the geometric mean $\sqrt{\lmin\lmax}$. With $\mub=1-h$, the self-dual line is the coupling
\begin{equation}
\Jcond=\frac{2(\mub-1)}{L}=2\Jup,\qquad \Jup=\frac{\mub-1}{L}.
\label{eq:thresholds}
\end{equation}

Both thresholds have a direct meaning. Take $\mub>1$, a rule that by itself favors short flights. After a flight at the cutoff, $x=L$, the field is $h+JL=JL-(\mub-1)$, and it is positive, so that the next flight is more likely long than short, when $J>\Jup$. Above this coupling a run of long flights can sustain itself. Whether the walker spends most of its time in such runs is decided by the duality: below $\Jcond$ it lives among short flights, above it among long ones. The fraction $\Pi$ of flights longer than $\sqrt{\lmin\lmax}$ is therefore the natural order parameter, and
\begin{equation}
\Phi=\frac J\Jup=\frac{JL}{\mub-1}
\end{equation}
the natural control parameter, with $\Phi=1$ at the first threshold and $\Phi=2$ at the second.

For wide ranges the change becomes a sharp transition. Since $\ln Z(\theta)=L\max(\theta,0)+O(\ln L)$, at fixed $\alpha=\mub-1>0$ and $\beta=JL$ the stationary law obeys a large-deviation principle in $L$,
\begin{equation}
\pi(x)\asymp e^{-L\,F(x/L)},\qquad F(u)=\alpha u-\max(\beta u-\alpha,0),
\label{eq:freeenergy}
\end{equation}
with $u=x/L\in[0,1]$ [\figref{fig:overview}(e)]. The logarithmic range $L$ plays the role of system size and $F$ that of a free energy per unit size. Because both factors of $\pi$ are exponentials in $x$, $F$ is piecewise linear and its minima can only sit at the ends, $u=0$ (flights near $\lmin$) and $u=1$ (flights near $\lmax$). For $\beta<2\alpha$ the global minimum is $u=0$ and $\Pi\to0$: the walker is in a Pareto phase, with a power-law histogram. For $\beta>2\alpha$ the global minimum is $u=1$ and $\Pi\to1$: the flights condense at the cutoff. The transition, on the self-dual line, is first order, with $\Pi$ jumping from 0 to 1. Between the two thresholds, $\alpha<\beta<2\alpha$, the end $u=1$ is a local minimum: a run of flights near $\lmax$ is metastable, and it lasts for a number of flights that grows exponentially with $L$ (Appendix~\ref{app:largeL}); $J=\Jup$ is the spinodal. At finite $L$ the jump is smoothed over a width of order $1/L$, and the curves $\Pi(\Phi)$ for different ranges cross at the self-dual point $\Phi=2$, $\Pi=1/2$ [\figref{fig:theory}(b)], as observables do at first-order transitions in finite systems~\cite{BinderLandau1984}.

Both thresholds vanish as $\lmax/\lmin\to\infty$. This reflects how the coupling scales, not a pathology. The coupling acts through $JL$, the largest change of field that one flight can impose on the next, and it is $JL$ that must be held fixed as the range grows, much as the couplings of a mean-field magnet must scale with the inverse number of spins. At fixed $J$ a wider range makes the same coupling stronger, and on the half-line every positive coupling lies beyond condensation, which is another way of saying that the chain does not exist there. At fixed $\Phi$ the limit $L\to\infty$ is well defined, and it is the limit in which the transition becomes sharp. The cutoff is therefore a physical ingredient of a L\'evy walk with memory, not a regularization to be removed at the end. Real data span one to two and a half decades, where the transition is a broad crossover, but its location still organizes the phase diagram in which the walkers of Sec.~\ref{sec:data} sit [\figref{fig:overview}(c)].

\subsection{Memory dresses the exponent}\label{sec:dressing}

A histogram does not know in which order the flights came. An i.i.d.\ truncated-Pareto fit to stationary flights matches their mean log length and returns $\hat\mu=1-\Lambda^{-1}(\langle x\rangle_\pi)$. For $J>0$ this exponent is smaller than $\mub$ [\figref{fig:theory}(c)]: memory packs long flights into runs, the histogram is heavier-tailed than the rule that generates it, and a walker with memory looks more L\'evy than its one-step rule. On the self-dual line $\langle x\rangle_\pi=L/2$ and $\hat\mu=1$ exactly, for every range. Put differently, the local exponent $\mu_t=1-\theta_t$ fluctuates with the recent past, a superstatistics~\cite{BeckCohen2003} generated by the dynamics itself. The exponent of a histogram is thus not a property of the search rule alone; it also depends on the memory and on the range over which it is measured. In real walkers the difference exceeds half a unit (Sec.~\ref{sec:data}).

\subsection{A critical point for the collective speed}\label{sec:meanfield}

Groups raise a question that single walkers do not: can a delayed coupling between walkers order the group as a whole? Take $N$ units with a global delayed coupling and no own memory, that is, Eq.~\eqref{eq:selfsocial} with $a_\tau=0$, $s_i$ the mean of all other units, a common field $h_i=H$ and no time fields. The group average $m_t$ plays the role of a magnetization. It is the collective log-speed, close to $0$ when the whole group crawls and close to $L$ when it runs. Two parameters control it. One is the total coupling $J_{\rm tot}=\sum_\tau J_\tau$, the strength with which the past speed of the group feeds back on its present. The other is a bias. Exchanging slow and fast, $m\to L-m$, maps the model onto itself with $H\to-H-J_{\rm tot}L$, the duality of Sec.~\ref{sec:condensation} at the level of the group; its symmetric point is $H=-J_{\rm tot}L/2$, and
\begin{equation}
\tilde\theta=H+\frac{J_{\rm tot}L}2
\end{equation}
measures the distance from it, positive when the group is biased toward running.

For $N\to\infty$ the units are independent given the past, and $m_t$ follows the deterministic delayed map
\begin{equation}
m_t=\Lambda\Big(H+\sum_\tau J_\tau m_{t-\tau}\Big),
\label{eq:map}
\end{equation}
with $\Lambda$ the Langevin response~\eqref{eq:langevin}. A steady state $m^*$, with field $\theta^*$, loses stability when a root $z$ of $1=\Lambda'(\theta^*)\sum_\tau J_\tau z^{-\tau}$ crosses the unit circle (Appendix~\ref{app:collective}). For positive couplings the first root to cross is $z=1$, a fold at which a second stable state appears. Since the susceptibility of one flight never exceeds $L^2/12$ (Sec.~\ref{sec:single}), this requires
\begin{equation}
J_{\rm tot}>J_c=\frac{12}{L^2},
\label{eq:Jc}
\end{equation}
with a cusp at $J_{\rm tot}=J_c$, $\tilde\theta=0$ [\figref{fig:many}(c)]. Above $J_c$ the group switches at random between a slow and a fast state, a collective stop-and-go whose residence times grow with $N$ [\figref{fig:many}(d)]. A negative coupling at delay $d$ makes the group oscillate instead: the map splits into $d$ independent sub-chains, each of which undergoes a period doubling at $|J|=J_c$, and any lag-1 inertia puts them in phase, so that the period is $2d$. Near both thresholds the fluctuations of the group average diverge, $N\,\mathrm{Var}(m)\to\Lambda'/[1-(\Lambda'J)^2]$ for one lag.

Own memory adds to the social feedback. With a lag-1 self coupling $a\ge0$, each unit at a steady state is the chain of Eq.~\eqref{eq:stationary} in the field $H+J_{\rm tot}m^*$, and the steady state loses stability when $J_{\rm tot}\chi=1$, where $\chi=\sum_{s\ge0}R(s)$ is the integrated response of one unit to a pulse of field. What enters is the response, not the integrated autocorrelation, which is 50\% larger at $a=0.2$. For $L=\ln100$ and $a=0.2$ the threshold drops to $J_c(a)=0.71J_c$, against $0.65J_c$ from the naive estimate $\Lambda'(0)(a+J)=1$, and the theory reproduces the simulated fluctuations within 1\% away from the threshold and within 4\% at $J=0.6J_c$ [Appendix~\ref{app:collective}, \figref{fig:collective}(b)].

The threshold $J_c=12/L^2$ exists because the support of each speed is compact, and it vanishes as the range grows. It is a scale-free and delayed counterpart of the speed criticality reported for starling flocks~\cite{Bialek2014}, where the interactions are instantaneous.
